\begin{table}[htbp]
\centering
\footnotesize
\caption{Problema D, meia-lua, Passos 2 a 5 (google/gemini-2.5-flash-lite): métrica em R$^2$, R$^3$ e R$^4$ por regime.}
\label{tab:D_passos_2_a_5_flashlite}
\begin{tabular}{lrrrrlrrlr}
\toprule
 &  &  & \multicolumn{3}{c}{Perceptron} & \multicolumn{3}{c}{NNLS} &  \\ \cmidrule(lr){4-6} \cmidrule(lr){7-9}
Regime & Acur.\ LLM vs real & Partição LLM (n0/n1) & Fid.\ (vs LLM) & Acur.\ vs real & $\hat{W}$ unit. & Fid.\ (vs LLM) & Acur.\ vs real & $\hat{W}$ unit. & $\Delta$ fid.\ vs R$^2$ \\
\midrule
R$^2$ — zero-shot† & 55,6 (53,3–59,3) & 86/64 & 81,8 (80,0–84,0) & 65,3 (64,7–66,0) & [0,39; 0,92] & 80,9 (80,0–82,0) & 66,2 (64,7–68,0) & [0,37; 0,93] & +0,0 \\
R$^2$ — 4-shot† & 77,9 (63,0–88,4) & 72/74 & 88,3 (72,6–97,9) & 77,9 (68,0–86,7) & [0,57; 0,80] & 87,9 (75,3–95,9) & 81,3 (70,0–86,7) & [0,28; 0,94] & +0,0 \\
R$^2$ — 10-shot† & 86,3 (81,4–90,0) & 71/69 & 95,2 (90,7–97,9) & 83,9 (81,3–86,7) & [0,35; 0,93] & 95,2 (90,7–98,6) & 84,4 (82,7–86,7) & [0,28; 0,96] & +0,0 \\
R$^2$ — 20-shot† & 88,5 (83,8–93,1) & 69/61 & 95,7 (94,6–97,7) & 85,7 (83,3–86,7) & [0,23; 0,97] & 95,0 (93,1–96,9) & 85,5 (83,3–87,3) & [0,25; 0,97] & +0,0 \\
R$^2$ — 40-shot† & 90,5 (86,4–93,6) & 56/54 & 94,7 (93,6–96,4) & 86,1 (84,0–88,0) & [0,25; 0,97] & 93,8 (90,9–95,5) & 85,7 (84,7–87,3) & [0,25; 0,97] & +0,0 \\
R$^2$ — W ajustado ao rótulo real (teto)† & — & — & — & 86,9 (86,7–87,3) & [0,22; 0,97] & — & 86,7 & [0,22; 0,97] & — \\
\midrule
R$^3$ — zero-shot† & 55,6 (53,3–59,3) & 86/64 & 85,3 (84,0–86,0) & 60,4 (57,3–62,7) & [0,42; 0,64; 0,64] & 84,4 (83,3–85,3) & 64,0 (61,3–66,0) & [0,24; 0,58; 0,75] & +3,6 \\
R$^3$ — 4-shot† & 77,9 (63,0–88,4) & 72/74 & 91,2 (84,2–97,9) & 78,0 (66,0–85,3) & [0,40; 0,62; 0,66] & 90,9 (80,8–96,6) & 79,8 (64,7–86,7) & [0,10; 0,54; 0,66] & +3,0 \\
R$^3$ — 10-shot† & 86,3 (81,4–90,0) & 71/69 & 95,8 (90,7–99,3) & 84,0 (81,3–86,0) & [0,26; 0,72; 0,63] & 95,4 (90,7–98,6) & 84,4 (82,0–86,7) & [0,19; 0,64; 0,49] & +0,6 \\
R$^3$ — 20-shot† & 88,5 (83,8–93,1) & 69/61 & 95,7 (93,8–97,7) & 85,0 (82,7–86,7) & [0,21; 0,73; 0,65] & 94,8 (93,1–97,7) & 84,5 (80,7–86,7) & [0,21; 0,70; 0,57] & +0,0 \\
R$^3$ — 40-shot† & 90,5 (86,4–93,6) & 56/54 & 95,2 (93,6–96,4) & 86,4 (84,7–87,3) & [0,21; 0,74; 0,64] & 94,0 (91,8–95,5) & 85,3 (82,7–87,3) & [0,23; 0,87; 0,22] & +0,4 \\
R$^3$ — W ajustado ao rótulo real (teto)† & — & — & — & 87,1 (86,7–87,3) & [0,16; 0,69; 0,71] & — & 86,7 & [0,22; 0,97; 0,00] & — \\
\midrule
R$^4$ — zero-shot† & 55,6 (53,3–59,3) & 86/64 & 79,8 (78,7–80,7) & 66,4 (64,7–68,7) & [0,24; 0,69; 0,10; 0,66] & 83,8 (80,7–88,7) & 62,4 (53,3–68,0) & [0,27; 0,63; 0,00; 0,33] & $-$2,0 \\
R$^4$ — 4-shot† & 77,9 (63,0–88,4) & 72/74 & 93,4 (87,7–97,3) & 78,4 (75,3–84,0) & [0,36; 0,63; 0,16; 0,58] & 92,5 (85,6–96,6) & 81,3 (76,7–85,3) & [0,30; 0,83; 0,13; 0,20] & +5,1 \\
R$^4$ — 10-shot† & 86,3 (81,4–90,0) & 71/69 & 96,2 (90,7–98,6) & 84,2 (82,0–86,0) & [0,22; 0,75; 0,07; 0,62] & 95,3 (89,3–98,6) & 84,1 (82,7–86,0) & [0,26; 0,93; 0,02; 0,12] & +1,0 \\
R$^4$ — 20-shot† & 88,5 (83,8–93,1) & 69/61 & 95,6 (93,8–96,9) & 85,7 (83,3–88,0) & [0,21; 0,75; 0,05; 0,62] & 95,0 (93,8–96,9) & 85,5 (83,3–87,3) & [0,23; 0,97; 0,02; 0,00] & $-$0,2 \\
R$^4$ — 40-shot† & 90,5 (86,4–93,6) & 56/54 & 95,5 (91,8–97,3) & 86,8 (85,3–88,7) & [0,16; 0,75; 0,04; 0,64] & 94,0 (90,9–96,4) & 85,9 (84,7–87,3) & [0,23; 0,96; 0,02; 0,07] & +0,7 \\
R$^4$ — W ajustado ao rótulo real (teto)† & — & — & — & 87,3 (86,7–88,7) & [0,15; 0,71; 0,03; 0,68] & — & 86,4 (86,0–86,7) & [0,22; 0,98; 0,00; 0,00] & — \\
\bottomrule
\multicolumn{10}{p{0.95\linewidth}}{\footnotesize Meia-lua (150 pontos, sem split). O LLM vê sempre só x1, x2; R$^3$ acrescenta x1·x2 e R$^4$ acrescenta x1², x2² apenas na estimação da métrica. Colunas e convenções da bloco1\_table\_01: o teto (W ajustado ao rótulo real) só tem acurácia (vs rótulo real), não fidelidade. Passos do Problema D: zero-shot = Passo 2 (o LLM rotula sem exemplos); $k$-shot = Passo 3 (o LLM rotula com $k$ exemplos do rótulo real); a métrica de cada linha é estimada no Passo 4; teto = Passo 5 (W ajustado diretamente ao rótulo real). Em zero-shot (e com nomes de classe neutros) a correspondência entre nome e classe é arbitrária: a acurácia (vs rótulo real) do LLM pode sair abaixo de 50\%; o valor invariante à permutação é max($a$, 1 $-$ $a$). † = em pelo menos uma estimativa o Perceptron não encontrou $\hat{W}$ que separe os rótulos: melhor esforço ($\gamma = 0$) ou margem bruta mínima negativa (algum ponto do lado errado da fronteira). ``—'' = colapso em classe única (ou quase-colapso: classe minoritária com menos de 5 pontos); na linha colapsada também a acurácia do LLM e a partição ficam como ``—''.} \\
\end{tabular}
\end{table}
